\documentclass{bmstu}

\begin{document}

\makereporttitle
{Информатика и системы управления (ИУ)}
{Программное обеспечение ЭВМ и информационные технологии (ИУ7)}
{\textbf{5}}
{Проектирование экспертных систем}
{Прямая дедукция для определенных выражений методом поиска в ширину}
{}
{ИУ7-32М}
{К.Э. Ковалец}
{З.Н. Русакова}


\setcounter{page}{2}
\renewcommand{\contentsname}{Содержание} 
\tableofcontents


\chapter{Основная часть}

\section{Унификация}

Унификация --- это процесс замены предметных переменных аргумента атома другими предметными переменными или константами или функциональными символами. Подстановка называется унификацией, если атомы совпадают.

Унификация выполняется:
\begin{itemize}
    \item если термы --- константы, то они унифицируются, если совпадают;
    \item если в первом атоме терм --- переменная, во втором --- константа, то они унифицируются, и переменная получает значение константы;
    \item если терм в первом атоме --- переменная и во втором --- переменная, то они унифицируются и становятся связанными, можно использовать любое имя.
\end{itemize}

Унифицируются между собой термы, стоящие на одинаковых местах в контрарных атомах.

\section{Доказательство правила}

В рамках процесса доказательства правила последовательно рассматриваются входные атомы или входные вершины текущего правила. Чтобы доказать входную вершину нужно найти в базе фактов такой атом, с кроторым будут совпадать имя рассматриваемого предиката и число аргументов, далее необходимо выполнить унификацию пары атомов. Если такой факт найден и унификация прошла успешна, то переменные рассматриваемого атома получают свои значения. Далее значения этих перемнных распространяются на остальные атомы данного правила, в которые они входят. После этого переходим к доказательству следующей входной вершины.

При доказательстве атома возможны два варианта: нужный факт в базе присутствует или отсутствует. Если факта нет, то входная вершина, а следовательно и само правило не будут доказаны. Если факт найден и унификация атомов прошла успешна, то полученные значения переменных распространяются по правилу. Если входные вершины доказаны, то доказан и выходной атом, который записывается в базу фактов или закрытые вершины. Правило считается доказанным и больше не рассматривается.

\section{Алгоритм прямого вывода}

Система дедукции на основе обобщённых правил продукции представляет собой алгоритм, который позволяет делать логические выводы на основе определённых выражений. В основе системы лежит метод резолюции, который позволяет обрабатывать более сложные представления знаний.

Алгоритм можно описать следующим образом:
\begin{enumerate}
    \item В список закрытых вершин (рабочую память) помещаются атомы, которые являются фактами.
    \item Передается целевой атом (возможно, с переменной, возможно, с константой), вводим два флага (<<решение найдено>>, <<решение не найдено>>).
    \item Пока оба флага равны $false$:
    \begin{itemize}
        \item запускается цикл по открытым правилам:
        \begin{itemize}
            \item проверяется возможность доказательсва правила;
            \item если оно доказано, то есть входные атомы правила входят в множество фактов, то:
            \begin{itemize}
                \item выходной атом правила, переменные в котором уже получили значения, добавляется к фактам (закрытым вершинам);
                \item доказанное правило делается закрытым.
            \end{itemize}
        \end{itemize}
    \end{itemize}
\end{enumerate}

\section{Пример работы программы}

Результаты работы программы приведены в листинге \ref{lst:result}.

\mylisting[text]{result.txt}{firstline=1,lastline=112}{Результат работы программы}{result}{}

\chapter{Текст программы}

В листингах \ref{lst:main}--\ref{lst:utils} представлен код программы.

\mylisting[python]{main.py}{firstline=1,lastline=37}{Основной модуль программы}{main}{}

\mylisting[python]{items.py}{firstline=1,lastline=102}{Модуль с описанием классов для термов, атомов, конъюнктов, импликаций и знаний}{items}{}

\mylisting[python]{unification.py}{firstline=1,lastline=64}{Модуль с реализацией функций унификации и подстановок}{unification}{}

\mylisting[python]{deduction.py}{firstline=1,lastline=118}{Модуль с реализацией класса прямой дедукции}{deduction}{}

\mylisting[python]{utils.py}{firstline=1,lastline=77}{Модуль с реализацией функций парсеров}{utils}{}

\end{document}
